% Supplementary material for the Physiological Measurement submission (compiled separately).
\documentclass{iopjournal}
\usepackage{amssymb}
\usepackage{amsmath}
\usepackage{booktabs}
\usepackage{array}
\usepackage{adjustbox}
\usepackage{makecell}
\usepackage{placeins}
\hypersetup{hidelinks}
\makeatletter
\setlength{\@fptop}{0pt}
\makeatother
\newcommand{\rev}[1]{#1}
% justified text (the class sets \raggedright globally)
\usepackage{ragged2e}
\AtBeginDocument{\justifying}
% Revision marking (2026-09-23 editorial and statistical revision): revised text is shown in blue.
\newif\ifshowchanges
\showchangestrue
\ifdefined\cleanbuild \showchangesfalse \fi
\newcommand{\chg}[1]{\ifshowchanges{\color{blue}#1}\else{#1}\fi}
\input{generated/results_macros.tex}
\fancyhead[L]{\small Supplementary material}
\fancyhead[C]{}
\fancyhead[R]{\small Subject-independent surface EMG electrode reduction}
\renewcommand{\thesection}{S\arabic{section}}
\renewcommand{\thetable}{S\arabic{table}}
\renewcommand{\thefigure}{S\arabic{figure}}

\begin{document}

{\noindent\fontsize{16}{19}\selectfont\sffamily Supplementary material\par}
\vspace{2mm}
{\noindent\large Subject-independent surface EMG electrode reduction using cross-fitted attribution: a prespecified non-inferiority study\par}
\vspace{6mm}

\chg{Supplementary methods S1--S6 are followed by supplementary figures S1--S2 and supplementary tables S1--S10, numbered in the order in which they are cited in the main text. Main-text tables are referred to as main table~1--3.}

\section{Departures from the analysis plan}
\label{app:deviations}
The following changes were made after the statistical analysis plan was written. None altered a prespecified decision rule (margin, statistical unit, bootstrap or definition of $k^*$).
\begin{enumerate}
\item The plan required Holm correction within H1. It was applied after the primary analysis and left $k^*$ unchanged at all three margins (supplementary table~S3). A fixed-sequence analysis, testing $k = 15, 14, \ldots$ and stopping at the first failure, is reported in the same table as a sensitivity analysis.
\item Training-sample sizes (40,000 windows for the support vector machine, 200,000 for the random forest and for the electrode-selection analyses) were chosen from held-out Macro-F1 on two folds (participants 1 and 7) before the main analyses. For these two participants the held-out data therefore informed a design choice.
\item The plan did not specify a criterion for choosing $\lambda$, so none was chosen and all four values are reported. The value $\lambda = 0.5$ used for SR-CFA in the H2 comparison was not chosen by a stated criterion, so the $\lambda = 0$ ranking is reported alongside it.
\item \chg{The plan named H1--H4 as confirmatory hypotheses with Holm correction within each family. The detailed H2--H4 test families (tests, rankings compared and subset sizes) were specified after the performance-curve and selection-method results were available; H2--H4 are therefore reported as secondary tests.}
\item The comparison model for H2 and H4 was reduced from 150 to 50 trees for computational cost before any comparison result was available.
\item An initial count-feature threshold 87 times smaller than the final GRABMyo value saturated WAMP and was replaced. The leave-one-out check of the thresholds (section~\ref{app:threshold}) was added after the primary analysis.
\item An initial deletion analysis combined deletion orders computed with different training sizes; it was replaced by the analysis reported here, in which all orders share one configuration (section~\ref{app:faith}).
\item The bipolar analysis was run three times. The first run used the monopolar count-feature threshold; the second used $5 \times 10^{-3}$~mV, whose stated derivation could not be reproduced; the reported analysis uses $4 \times 10^{-3}$~mV, the value given by the documented rule. Earlier outputs are archived.
\item The fixed-layout and feature-redundancy analyses (sections~\ref{app:fixed} and \ref{app:redund}) were not in the plan. \chg{The fixed-layout check used $\lambda = 0.5$; no rationale for this value was recorded (section~\ref{app:fixed}).}
\item \chg{The participant-level attribution extension (section~\ref{app:participant_extension}) was specified after the original results were available.}
\item The window-length sensitivity analysis and the 17-class task, listed as exploratory in the plan, were not run.
\end{enumerate}

\section{Leave-one-out check of the count-feature thresholds}
\label{app:threshold}
All thresholds follow one rule: 1.15 times the median absolute first difference of the signal, rounded to one significant figure. The GRABMyo threshold was derived from 54 unfiltered recordings of 6 participants. For every outer fold we recomputed the median without the held-out participant and applied the same rule. It gave $10^{-2}$~mV in \LOOMonoIdentical\ of \LOOMonoFolds\ folds (leave-one-out medians \LOOMonoMin--\LOOMonoMax${}\times 10^{-3}$~mV); only \LOOMonoInProbe\ folds held out a participant who contributed to the sample. For the derived bipolar signal (median \LOOBipPooled${}\times 10^{-3}$~mV on the same recordings) it gave $\LOOBipRule \times 10^{-3}$~mV in \LOOBipIdentical\ of \LOOBipFolds\ folds. For DB5 it gave 7 quantisation steps in \LOODBFiveIdentical\ of \LOODBFiveFolds\ folds. Deriving the thresholds within each fold would therefore have produced identical features. Because the thresholds were not estimated within folds, we describe this as a verification rather than as fold-local estimation.

\section{Consistency of the deletion analysis}
\label{app:faith}
The Remove-Bottom subsets coincide with the H2 cross-fitted SHAP subsets in all \FaithProtocolCells\ participant-by-size cells, and their Macro-F1 values agree to within a mean absolute difference of \FaithProtocolMeanDiff\ (maximum \FaithProtocolMaxDiff). The small differences arise because features were standardised separately for each subset in the deletion analysis and once over all features in the H2 comparison.

\section{Fixed-layout consistency check}
\label{app:fixed}
A single 12-electrode layout was derived from a cross-fitted ranking on the \FixedDev\ odd-numbered participants, removing electrodes \FixedDropped. Models with 16 and with 12 electrodes were trained on those participants and tested on the \FixedTest\ even-numbered participants. The 12-electrode model scored $\FixedDiff$ against the 16-electrode model (95\% confidence interval $\FixedLCB$ to $\FixedUCB$), within the margin. The electrode count was chosen after inspecting the full-cohort performance curve, which included the even-numbered participants. Although electrode identities were selected using development participants only, this is therefore not an independent confirmation of the chosen count. It is a single-split consistency check \chg{and does not constitute a prospective validation of a fixed sparse hardware layout.}

\chg{The ranking for this check used $\lambda = 0.5$, whereas the primary analysis uses $\lambda = 0$; no rationale for this choice was recorded. In the primary LOSO analysis the $\lambda = 0.5$ and $\lambda = 0$ rankings selected a different top-12 set in \TopTwelveDiffersLambdaPtFive\ of 43 outer folds, but this was not verified for the development split used here. The removed set differs by one electrode from the set excluded in the primary outer folds (\DroppedFourSet).} [AUTHOR CHECK REQUIRED: justify the use of $\lambda = 0.5$ in the fixed-layout check, or repeat the check with $\lambda = 0$.]

\section{Feature redundancy}
\label{app:redund}
Integrated EMG equals 512 times the mean absolute value for these windows. Repeating the comparison-model evaluation without integrated EMG changed Macro-F1 by $\RedundDiff$ (95\% confidence interval $\RedundLCB$ to $\RedundUCB$; 43 folds). The median per-fold Spearman $\rho$ between rankings with and without integrated EMG was \RedundRho, and the top-12 set was identical in \RedundTopTwelveSame\ of 43 folds, so the duplicated feature had little measurable effect on the ranking.

\section{Participant-level attribution extension}
\label{app:participant_extension}
The post-hoc extension uses the same out-of-fold attributions as the original ranking, but gives each outer-training participant equal weight before computing the median and median absolute deviation. All penalties and subset sizes were evaluated with the same training budget as the original curve. Every reused fit records its source, and inputs are fingerprinted to prevent mixing incompatible runs. Between-person dispersion is descriptive heterogeneity, not a standard error of attribution. It includes both participant differences and differences between the inner fitted models.

\input{generated/participant_extension_results.tex}

\begin{figure}[htbp]
\centering
\includegraphics[width=\linewidth]{figures/fig9_participant_extension.pdf}
\caption{\chg{Participant-level versus fold-level attribution aggregation (post hoc; GRABMyo, 43 participants).} Every cell is the mean paired Macro-F1 difference across participants, comparing person-level with fold-level aggregation under the same penalty, subset size and model protocol. Blue favours person-level aggregation and orange favours fold aggregation. Asterisks identify Holm-corrected differences within the complete family. All candidates are shown; outer-test outcomes did not select a penalty.}
\label{fig:participant_extension}
\end{figure}

\begin{figure}[htbp]
\centering
\includegraphics[width=\linewidth]{figures/fig10_penalty_gate.pdf}
\caption{\chg{Stability penalty versus attribution magnitude alone under participant-level aggregation (post hoc).} (a) $\lambda = 0.25$, (b) $\lambda = 0.5$, (c) $\lambda = 1$. Each point compares a positive penalty with magnitude alone. Bars are unadjusted participant-bootstrap 95\% confidence intervals; asterisks indicate Holm-adjusted $p<0.05$ across all penalty-by-size comparisons. Positive differences favour the penalty. Zero-width intervals arise when all paired results coincide. This diagnostic is exploratory and does not replace the original H3 test.}
\label{fig:penalty_gate}
\end{figure}

\FloatBarrier
\section*{Supplementary tables}
\begin{table*}[htbp]
\centering
\small
\caption{\chg{Full-array (16-electrode) leave-one-subject-out performance of three classifiers on GRABMyo (43 participants, 16 active gestures; chance Macro-F1 $\approx$ 0.063). CI: 95\% participant cluster-bootstrap confidence interval. RBF-SVM: support vector machine with a radial basis function kernel. Training windows: windows used to fit each fold's model.}}
\label{tab:full_channel}
\setlength{\tabcolsep}{4pt}
\begin{adjustbox}{max width=\textwidth}
\begin{tabular}{llcccc}
\toprule
Model & Training windows & Macro-F1 & 95\% CI & SD & Range \\
\midrule
XGBoost & \chg{all (550{,}368 per fold)} & 0.542 & 0.517--0.567 & 0.085 & 0.385--0.702 \\
RBF-SVM & 40{,}000 & 0.538 & 0.513--0.562 & 0.085 & 0.383--0.718 \\
Random forest & 200{,}000 & 0.502 & 0.475--0.528 & 0.090 & 0.346--0.670 \\
\bottomrule
\end{tabular}
\end{adjustbox}
\par\smallskip\parbox{0.95\linewidth}{\footnotesize Paired comparisons (Wilcoxon signed-rank, Holm over 3): Random forest $-$ RBF-SVM = $-$0.0362 (95\% CI $-$0.0444 to $-$0.0284), $p_{\mathrm{Holm}}$ $<$ 0.001, $r_{rb}$ = $-$0.97; Random forest $-$ XGBoost = $-$0.0408 (95\% CI $-$0.0475 to $-$0.0340), $p_{\mathrm{Holm}}$ $<$ 0.001, $r_{rb}$ = $-$0.98; RBF-SVM $-$ XGBoost = $-$0.0045 (95\% CI $-$0.0107 to +0.0015), $p_{\mathrm{Holm}}$ = 0.154, $r_{rb}$ = $-$0.25.}
\end{table*}

\begin{table}[htbp]
\centering
\small
\caption{\chg{Full-array XGBoost F1 score by active GRABMyo gesture (43 participants). Mean F1 across participants with a 95\% percentile cluster-bootstrap confidence interval of the mean (10,000 resamples). Gesture names follow the order given in the dataset file MotionSequence.txt.}}
\label{tab:gesture_audit}
\begin{tabular}{rlrr}
\toprule
ID & \chg{Gesture} & Mean F1 & \chg{95\% bootstrap CI} \\
\midrule
1 & Lateral prehension & 0.444 & 0.389 to 0.496 \\
2 & Thumb adduction & 0.309 & 0.256 to 0.360 \\
3 & Thumb and little finger opposition & 0.526 & 0.472 to 0.577 \\
4 & Thumb and index finger opposition & 0.349 & 0.297 to 0.399 \\
5 & Thumb and index finger extension & 0.486 & 0.424 to 0.547 \\
6 & Thumb and little finger extension & 0.602 & 0.532 to 0.667 \\
7 & Index and middle finger extension & 0.506 & 0.443 to 0.566 \\
8 & Little finger extension & 0.499 & 0.441 to 0.553 \\
9 & Index finger extension & 0.388 & 0.333 to 0.443 \\
10 & Thumb finger extension & 0.538 & 0.470 to 0.603 \\
11 & Wrist extension & 0.942 & 0.927 to 0.956 \\
12 & Wrist flexion & 0.828 & 0.775 to 0.874 \\
13 & Forearm supination & 0.814 & 0.770 to 0.852 \\
14 & Forearm pronation & 0.491 & 0.431 to 0.550 \\
15 & Hand open & 0.545 & 0.489 to 0.598 \\
16 & Hand close & 0.410 & 0.368 to 0.449 \\
\bottomrule
\end{tabular}
\end{table}

\FloatBarrier
\begin{table*}[htbp]
\centering
\small
\caption{\chg{H1 decisions under multiplicity correction and sensitivity margins (GRABMyo, 43 participants, $\lambda = 0$). $\bar\Delta_k$ and lower 95\% bound as in main table~1. $p$: one-sided cluster-bootstrap $p$ for $H_0\colon \bar\Delta_k \le -\delta$ at the primary margin $\delta = 0.02$ (10,000 resamples); Holm over the 15 subset sizes at $\alpha = 0.025$. Fixed sequence (post-hoc sensitivity analysis): $k = 15, 14, \ldots$ tested in order, stopping at the first failure. The last two columns give the non-inferiority decision (unadjusted / Holm) at the prespecified sensitivity margins.}}
\label{tab:h1_holm}
\setlength{\tabcolsep}{4pt}
\begin{adjustbox}{max width=\textwidth}
\begin{tabular}{rcccccccccc}
\toprule
$k$ & $\bar\Delta_k$ & Lower bound & $p$ & $p_{\mathrm{Holm}}$ & Non-inf. & Non-inf. (Holm) & Fixed sequence & $\delta = 0.01$ & $\delta = 0.03$ \\
\midrule
15 & $-$0.0036 & $-$0.0056 & $<$0.001 & 0.002 & yes & yes & yes & yes / yes & yes / yes \\
14 & $-$0.0060 & $-$0.0081 & $<$0.001 & 0.002 & yes & yes & yes & yes / yes & yes / yes \\
13 & $-$0.0109 & $-$0.0133 & $<$0.001 & 0.002 & yes & yes & yes & no / no & yes / yes \\
12 & $-$0.0133 & $-$0.0160 & $<$0.001 & 0.002 & yes & yes & yes & no / no & yes / yes \\
11 & $-$0.0204 & $-$0.0246 & 0.579 & 1.000 & no & no & no & no / no & yes / yes \\
10 & $-$0.0259 & $-$0.0311 & 0.984 & 1.000 & no & no & no & no / no & no / no \\
9 & $-$0.0318 & $-$0.0387 & 0.999 & 1.000 & no & no & no & no / no & no / no \\
8 & $-$0.0366 & $-$0.0440 & 1.000 & 1.000 & no & no & no & no / no & no / no \\
7 & $-$0.0413 & $-$0.0489 & 1.000 & 1.000 & no & no & no & no / no & no / no \\
6 & $-$0.0610 & $-$0.0713 & 1.000 & 1.000 & no & no & no & no / no & no / no \\
5 & $-$0.1133 & $-$0.1277 & 1.000 & 1.000 & no & no & no & no / no & no / no \\
4 & $-$0.1444 & $-$0.1590 & 1.000 & 1.000 & no & no & no & no / no & no / no \\
3 & $-$0.2015 & $-$0.2212 & 1.000 & 1.000 & no & no & no & no / no & no / no \\
2 & $-$0.2420 & $-$0.2603 & 1.000 & 1.000 & no & no & no & no / no & no / no \\
1 & $-$0.3603 & $-$0.3832 & 1.000 & 1.000 & no & no & no & no / no & no / no \\
\bottomrule
\end{tabular}
\end{adjustbox}
\par\smallskip\parbox{0.95\linewidth}{\footnotesize $k^*$ (unadjusted / Holm / fixed sequence) at $\lambda = 0$: $\delta$ = 0.01: 14 / 14 / 14; $\delta$ = 0.02: 12 / 12 / 12; $\delta$ = 0.03: 11 / 11 / 11; identical for all four $\lambda$.}
\end{table*}

\begin{table}[htbp]
\centering
\small
\caption{\chg{Exploratory comparisons of SR-CFA ($\lambda = 0.5$) with other electrode-selection methods (GRABMyo, 43 participants; same comparison model as main table~2). Difference: SR-CFA minus comparator, mean paired difference with 95\% participant cluster-bootstrap confidence interval (CI). $p$: Wilcoxon signed-rank, uncorrected; $r_{rb}$: matched-pairs rank-biserial correlation. These contrasts are exploratory and were not corrected for multiplicity.}}
\label{tab:selection_explor}
\setlength{\tabcolsep}{4pt}
\begin{adjustbox}{max width=\textwidth}
\begin{tabular}{rlcccccc}
\toprule
$k$ & Comparator & SR-CFA & Comparator & Difference & 95\% CI & $p$ & $r_{rb}$ \\
\midrule
4 & Permutation importance & 0.383 & 0.385 & $-$0.0018 & $-$0.0151 to +0.0118 & 0.678 & $-$0.08 \\
 & Geometric farthest-point & 0.383 & 0.383 & $-$0.0005 & $-$0.0159 to +0.0143 & 0.986 & +0.00 \\
 & Mutual information & 0.383 & 0.280 & +0.1031 & +0.0880 to +0.1176 & $<$0.001 & +0.99 \\
\midrule 6 & Permutation importance & 0.453 & 0.440 & +0.0126 & +0.0047 to +0.0216 & 0.009 & +0.52 \\
 & Geometric farthest-point & 0.453 & 0.419 & +0.0341 & +0.0195 to +0.0476 & $<$0.001 & +0.68 \\
 & Mutual information & 0.453 & 0.360 & +0.0934 & +0.0813 to +0.1050 & $<$0.001 & +1.00 \\
\midrule 8 & Permutation importance & 0.469 & 0.466 & +0.0033 & $-$0.0006 to +0.0071 & 0.069 & +0.33 \\
 & Geometric farthest-point & 0.469 & 0.465 & +0.0047 & $-$0.0044 to +0.0137 & 0.267 & +0.20 \\
 & Mutual information & 0.469 & 0.396 & +0.0737 & +0.0632 to +0.0841 & $<$0.001 & +1.00 \\
\midrule 12 & Permutation importance & 0.485 & 0.485 & +0.0001 & $-$0.0033 to +0.0033 & 0.995 & $-$0.00 \\
 & Geometric farthest-point & 0.485 & 0.478 & +0.0070 & +0.0000 to +0.0138 & 0.058 & +0.33 \\
 & Mutual information & 0.485 & 0.480 & +0.0051 & +0.0013 to +0.0089 & 0.007 & +0.47 \\
\bottomrule
\end{tabular}
\end{adjustbox}
\end{table}

\FloatBarrier
\begin{table}[htbp]
\centering
\small
\caption{\chg{H4: deletion with retraining (GRABMyo, 43 participants). Retained: electrodes kept after deletion. Remove-Top, Remove-Bottom and Remove-Random delete the highest-ranked, lowest-ranked or randomly chosen electrodes; Remove-Random uses the random subsets of the H2 comparison and exists only where these were sampled. All orders use the same comparison model and training sample. Difference: mean paired difference (first minus second order) with 95\% participant cluster-bootstrap confidence interval (CI); negative values follow the expected ordering. $p_{\mathrm{Holm}}$: Wilcoxon signed-rank, Holm over the 15 contrasts; $r_{rb}$: matched-pairs rank-biserial correlation.}}
\label{tab:faithfulness}
\setlength{\tabcolsep}{4pt}
\begin{adjustbox}{max width=\textwidth}
\begin{tabular}{rlcccccc}
\toprule
Retained & Contrast & First & Second & Difference & 95\% CI & $p_{\mathrm{Holm}}$ & $r_{rb}$ \\
\midrule
2 & Top $-$ Bottom & 0.263 & 0.292 & $-$0.0297 & $-$0.0419 to $-$0.0166 & $<$0.001 & $-$0.67 \\
\midrule 4 & Top $-$ Bottom & 0.351 & 0.382 & $-$0.0314 & $-$0.0447 to $-$0.0177 & $<$0.001 & $-$0.64 \\
 & Top $-$ Random & 0.351 & 0.366 & $-$0.0151 & $-$0.0239 to $-$0.0061 & 0.010 & $-$0.51 \\
 & Random $-$ Bottom & 0.366 & 0.382 & $-$0.0163 & $-$0.0265 to $-$0.0061 & 0.013 & $-$0.47 \\
\midrule 6 & Top $-$ Bottom & 0.376 & 0.454 & $-$0.0775 & $-$0.0921 to $-$0.0627 & $<$0.001 & $-$0.96 \\
 & Top $-$ Random & 0.376 & 0.416 & $-$0.0398 & $-$0.0487 to $-$0.0307 & $<$0.001 & $-$0.91 \\
 & Random $-$ Bottom & 0.416 & 0.454 & $-$0.0377 & $-$0.0459 to $-$0.0293 & $<$0.001 & $-$0.93 \\
\midrule 8 & Top $-$ Bottom & 0.408 & 0.469 & $-$0.0618 & $-$0.0772 to $-$0.0461 & $<$0.001 & $-$0.89 \\
 & Top $-$ Random & 0.408 & 0.445 & $-$0.0377 & $-$0.0476 to $-$0.0273 & $<$0.001 & $-$0.86 \\
 & Random $-$ Bottom & 0.445 & 0.469 & $-$0.0242 & $-$0.0309 to $-$0.0172 & $<$0.001 & $-$0.86 \\
\midrule 10 & Top $-$ Bottom & 0.445 & 0.475 & $-$0.0297 & $-$0.0399 to $-$0.0188 & $<$0.001 & $-$0.73 \\
\midrule 12 & Top $-$ Bottom & 0.455 & 0.485 & $-$0.0302 & $-$0.0385 to $-$0.0218 & $<$0.001 & $-$0.84 \\
 & Top $-$ Random & 0.455 & 0.478 & $-$0.0228 & $-$0.0296 to $-$0.0158 & $<$0.001 & $-$0.81 \\
 & Random $-$ Bottom & 0.478 & 0.485 & $-$0.0074 & $-$0.0101 to $-$0.0046 & $<$0.001 & $-$0.74 \\
\midrule 14 & Top $-$ Bottom & 0.487 & 0.491 & $-$0.0038 & $-$0.0084 to +0.0009 & 0.150 & $-$0.25 \\
\bottomrule
\end{tabular}
\end{adjustbox}
\end{table}

\begin{table}[htbp]
\centering
\small
\caption{\chg{H3: stability-regularised ranking ($\lambda > 0$) minus magnitude-only ranking ($\lambda = 0$) on the electrode-count performance curve (GRABMyo, 43 participants). Mean paired difference with 95\% participant cluster-bootstrap confidence interval (CI). Tied: participants for whom both rankings select the identical subset. $p_{\mathrm{Holm}}$: Wilcoxon signed-rank, Holm over 12 tests. $r_{rb}$ is computed from untied participants only and is undefined (--) when all are tied; with most participants tied it is not informative about magnitude.}}
\label{tab:lambda}
\setlength{\tabcolsep}{4pt}
\begin{adjustbox}{max width=\textwidth}
\begin{tabular}{rrccccc}
\toprule
$\lambda$ & $k$ & Mean difference & 95\% CI & Tied & $p_{\mathrm{Holm}}$ & $r_{rb}$ \\
\midrule
0.25 & 4 & $-$0.0002 & $-$0.0005 to +0.0000 & 41/43 & 1.000 & $-$1.00 \\
0.25 & 6 & +0.0000 & +0.0000 to +0.0000 & 43/43 & 1.000 & -- \\
0.25 & 8 & +0.0001 & +0.0000 to +0.0002 & 42/43 & 1.000 & +1.00 \\
0.25 & 12 & +0.0000 & +0.0000 to +0.0000 & 43/43 & 1.000 & -- \\
0.5 & 4 & +0.0004 & $-$0.0005 to +0.0018 & 40/43 & 1.000 & +0.00 \\
0.5 & 6 & +0.0000 & +0.0000 to +0.0000 & 43/43 & 1.000 & -- \\
0.5 & 8 & +0.0008 & +0.0000 to +0.0019 & 40/43 & 1.000 & +1.00 \\
0.5 & 12 & +0.0000 & +0.0000 to +0.0000 & 43/43 & 1.000 & -- \\
1 & 4 & +0.0010 & $-$0.0003 to +0.0027 & 38/43 & 1.000 & +0.47 \\
1 & 6 & +0.0000 & +0.0000 to +0.0000 & 43/43 & 1.000 & -- \\
1 & 8 & +0.0015 & +0.0004 to +0.0028 & 36/43 & 0.216 & +1.00 \\
1 & 12 & $-$0.0005 & $-$0.0014 to +0.0000 & 42/43 & 1.000 & $-$1.00 \\
\bottomrule
\end{tabular}
\end{adjustbox}
\end{table}

\FloatBarrier
\begin{table}[htbp]
\centering
\small
\caption{Participant-level attribution extension (exploratory; GRABMyo). Non-inferiority uses the same-budget full-array reference and the original engineering margin. The paired difference compares participant versus fold aggregation at the same penalty and electrode count.}
\label{tab:participant_extension}
\begin{tabular}{rrrrr}
\toprule
$\lambda$ & $k^*$ & $k^*$ (Holm) & $\Delta$ at $k=12$ (95\% CI) & $p_{\rm Holm}$ \\
\midrule
0 & 12 & 12 & +0.0000 (+0.0000, +0.0000) & 1.000 \\
0.25 & 12 & 12 & +0.0000 (+0.0000, +0.0000) & 1.000 \\
0.5 & 12 & 12 & +0.0000 (+0.0000, +0.0000) & 1.000 \\
1 & 12 & 12 & +0.0005 (+0.0000, +0.0014) & 1.000 \\
\bottomrule
\end{tabular}
\end{table}

\begin{table}[htbp]
\centering
\small
\caption{\chg{NinaPro DB5 channel-count performance curve (exploratory replication; cross-fitted ranking, $\lambda = 0$; 10 subjects; reference $k = 16$, mean Macro-F1 0.281; margin $\delta = 0.02$). The lower 95\% bound is shown to six decimals because the $k = 12$ decision lies at the margin. Non-inferior: lower bound above $-\delta$, unadjusted and after Holm correction over 15 subset sizes. Bootstrap reruns: number of 200 reruns of the bootstrap with different random seeds in which the unadjusted decision was non-inferior. This is a Monte Carlo decision-stability diagnostic, not a probability that the hypothesis is true.}}
\label{tab:replication}
\setlength{\tabcolsep}{4pt}
\begin{adjustbox}{max width=\textwidth}
\begin{tabular}{rcccccc}
\toprule
$k$ & Macro-F1 & $\Delta_k$ & Lower bound & Non-inf. & Non-inf. (Holm) & Non-inf. in bootstrap reruns \\
\midrule
15 & 0.270 & $-$0.0108 & $-$0.023773 & no & no & 0/200 (0.0\%) \\
14 & 0.269 & $-$0.0119 & $-$0.027438 & no & no & 0/200 (0.0\%) \\
13 & 0.266 & $-$0.0143 & $-$0.031186 & no & no & 0/200 (0.0\%) \\
12 & 0.268 & $-$0.0126 & $-$0.019994 & yes & no & 191/200 (95.5\%) \\
11 & 0.260 & $-$0.0211 & $-$0.034121 & no & no & 0/200 (0.0\%) \\
10 & 0.265 & $-$0.0160 & $-$0.031368 & no & no & 0/200 (0.0\%) \\
9 & 0.260 & $-$0.0202 & $-$0.036943 & no & no & 0/200 (0.0\%) \\
8 & 0.252 & $-$0.0288 & $-$0.055369 & no & no & 0/200 (0.0\%) \\
7 & 0.232 & $-$0.0486 & $-$0.075346 & no & no & 0/200 (0.0\%) \\
6 & 0.231 & $-$0.0498 & $-$0.081816 & no & no & 0/200 (0.0\%) \\
5 & 0.200 & $-$0.0803 & $-$0.118422 & no & no & 0/200 (0.0\%) \\
4 & 0.180 & $-$0.1006 & $-$0.150266 & no & no & 0/200 (0.0\%) \\
3 & 0.162 & $-$0.1185 & $-$0.174145 & no & no & 0/200 (0.0\%) \\
2 & 0.138 & $-$0.1430 & $-$0.195307 & no & no & 0/200 (0.0\%) \\
1 & 0.105 & $-$0.1754 & $-$0.229548 & no & no & 0/200 (0.0\%) \\
\bottomrule
\end{tabular}
\end{adjustbox}
\end{table}

\FloatBarrier
\begin{table}[htbp]
\centering
\small
\caption{Cross-session evaluation on GRABMyo (exploratory; full array). All conditions are also subject-independent: the test participant is never in training.}
\label{tab:cross_session}
\setlength{\tabcolsep}{4pt}
\begin{adjustbox}{max width=\textwidth}
\begin{tabular}{lcccc}
\toprule
Design & Train sessions & Test session & Train windows & Macro-F1 \\
\midrule
CS1 & 1 & 1 & 183{,}456 & 0.492 $\pm$ 0.101 \\
CS1 & 1 & 2 & 183{,}456 & 0.528 $\pm$ 0.088 \\
CS1 & 1 & 3 & 183{,}456 & 0.510 $\pm$ 0.090 \\
CS2 & 1+2 & 3 & 366{,}912 & 0.530 $\pm$ 0.095 \\
\bottomrule
\end{tabular}
\end{adjustbox}
\par\smallskip\parbox{0.95\linewidth}{\footnotesize Paired contrasts (exploratory, uncorrected): CS1: test S2 $-$ test S1: +0.0355 (95\% CI +0.0091 to +0.0621), $p$ = 0.024; CS1: test S3 $-$ test S1: +0.0172 (95\% CI $-$0.0038 to +0.0391), $p$ = 0.272; S3 test: CS2 (train S1+S2) $-$ CS1 (train S1): +0.0210 (95\% CI +0.0146 to +0.0274), $p$ $<$ 0.001.}
\end{table}

\begin{table}[htbp]
\centering
\small
\caption{\chg{GRABMyo derived bipolar pairs $F_j - F_{j+8}$ (exploratory; 43 participants; $\lambda = 0$; count-feature threshold $4 \times 10^{-3}$~mV; margin $\delta = 0.02$). Reference: the $k = 8$ model of the same curve, trained on the curve's 200,000-window sample (mean Macro-F1 0.480). The full-array bipolar value reported in the main text (0.478) comes from a separate model trained on all 550{,}368 training windows of each fold, as for the monopolar analysis. $\Delta_k$: mean paired difference with 95\% participant cluster-bootstrap confidence interval (CI). Bootstrap reruns: as in supplementary table~S8.}}
\label{tab:bipolar}
\setlength{\tabcolsep}{4pt}
\begin{adjustbox}{max width=\textwidth}
\begin{tabular}{rccccc}
\toprule
$k$ & Macro-F1 & $\Delta_k$ & 95\% CI & Non-inf. & Non-inf. in bootstrap reruns \\
\midrule
7 & 0.464 & $-$0.0158 & $-$0.0203 to $-$0.0113 & no & 0/200 (0.0\%) \\
6 & 0.442 & $-$0.0375 & $-$0.0451 to $-$0.0303 & no & 0/200 (0.0\%) \\
5 & 0.405 & $-$0.0742 & $-$0.0845 to $-$0.0646 & no & 0/200 (0.0\%) \\
4 & 0.356 & $-$0.1240 & $-$0.1392 to $-$0.1096 & no & 0/200 (0.0\%) \\
3 & 0.310 & $-$0.1697 & $-$0.1852 to $-$0.1546 & no & 0/200 (0.0\%) \\
2 & 0.240 & $-$0.2399 & $-$0.2615 to $-$0.2177 & no & 0/200 (0.0\%) \\
1 & 0.174 & $-$0.3052 & $-$0.3298 to $-$0.2790 & no & 0/200 (0.0\%) \\
\bottomrule
\end{tabular}
\end{adjustbox}
\end{table}


\end{document}
